% Generated draft: requires mathematical and source review.
\documentclass[11pt]{article}
\usepackage[margin=25mm]{geometry}
\usepackage{fontspec}
\setmainfont{Latin Modern Roman}
\usepackage{amsmath,amssymb,mathtools,unicode-math}
\setmathfont{Latin Modern Math}
\usepackage{enumitem,needspace,xurl,hyperref}
\setlength{\parindent}{0pt}
\setlength{\parskip}{5pt}
\setlength{\emergencystretch}{3em}
\newcommand{\N}{\mathbb{N}}
\newcommand{\Z}{\mathbb{Z}}
\newcommand{\Q}{\mathbb{Q}}
\newcommand{\R}{\mathbb{R}}
\newcommand{\C}{\mathbb{C}}
\newcommand{\F}{\mathbb{F}}
\newcommand{\A}{\mathbb{A}}
\newcommand{\PP}{\mathbb{P}}
\newcommand{\E}{\mathbb{E}}
\newcommand{\eps}{\varepsilon}
\newcommand{\dd}{\,\mathrm{d}}
\begin{document}
\section*{UnsolvedMath ID: 20002445}
\textbf{Source problem code:} catalogue entry
\subsection{Definitions and mathematical statement}
\textbf{Recovered Sandstede problem.} The question ``To what extent is this true for more general parabolic PDEs and for other nonlinear equations?'' refers, according to the supplied boundary recovery, to extension of Evans functions to branch points of the dispersion relation and to Murata's scalar-heat results relating temporal decay to roots at those points. The operator, boundary conditions, Evans normalization and decay norms are not specified; it is not a complete universal decay theorem.

\textbf{Supplied branch-to-time diagnostic.} Let a verified pointwise or weighted resolvent have a nonanalytic term $B(\lambda-\lambda_*)^\alpha$, using the branch cut $\lambda_*-[0,\infty)$ and the principal power with arguments $\pm\pi$ on its sides. Assume contour deformation, remainder bounds and spectral separation justify termwise long-time inversion. For a real noninteger $\alpha>-1$, the aid's sign-exact contribution is
\[
 e^{\lambda_*t}\,\frac{B}{\Gamma(-\alpha)}\,t^{-\alpha-1}.
\]
For data $f$ the coefficient is $Bf$ and vanishes if $Bf=0$; the next term must then be examined. The sign is not obtained by replacing $\Gamma(-\alpha)$ with its absolute value. A pole term $B(\lambda-\lambda_*)^{-m}$ contributes $e^{\lambda_*t}Bt^{m-1}/(m-1)!$. Nonnegative integer powers are locally analytic and by themselves give no branch tail for $t>0$; logarithmic terms require the corresponding parameter derivative/limiting inversion, not the generic power formula blindly. More singular noninteger exponents need a finite-part or continuation justification. A dispersion double root or Evans zero alone does not establish that a nonzero resolvent coefficient $B$ exists.

\textbf{Supplied exact half-line models.} For $u_t=u_{xx}$ on $x>0$, put $q=\sqrt\lambda$, $\operatorname{Re}q>0$. The resolvent kernels are
\[
 R_N(x,y;\lambda)=\frac{e^{-q|x-y|}+e^{-q(x+y)}}{2q},\qquad
 R_D(x,y;\lambda)=\frac{e^{-q|x-y|}-e^{-q(x+y)}}{2q}.
\]
A compatible uniformizer Evans normalization has $E_N(q)=q$ and $E_D(q)=1$. The Neumann zero at $q=0$ corresponds to the bounded constant threshold solution, which is not $L^2(0,\infty)$; it is not an eigenvalue claim. The Dirichlet singular $q^{-1}$ terms cancel. Their heat kernels for fixed $x,y>0$ satisfy
\[
 K_N(t,x,y)\sim\frac1{\sqrt\pi}\,t^{-1/2},\qquad
 K_D(t,x,y)\sim\frac{xy}{2\sqrt\pi}\,t^{-3/2}.
\]
The latter follows from the resolvent coefficient $-xy\sqrt\lambda$ and $\Gamma(-1/2)=-2\sqrt\pi$, displaying the sign. For suitably integrable data, for example compact support, the leading local Neumann coefficient depends on $\int f$ and the Dirichlet coefficient on $\int yf(y)\,\dd y$; either can cancel. These are fixed-point/local tails, not uniform $L^2$ semigroup norm decay. The paired models share the dispersion branch point but have different boundary/resolvent numerators. They give a scoped mechanism rather than a general systems or nonlinear stability theorem.

\textbf{Separate S.~Malham ``Biscale chaos'' setup.} The appended next-presenter material concerns
\[
 u_t=\delta\Delta u-uv^2,\qquad v_t=\Delta v+uv^2,
 \qquad\delta>0,
\]
a model of autocatalysis. The source describes $\delta$ as a diffusion/speed-ratio parameter, reports complicated perturbed fronts near $\delta\sim8$, and cites a suggestion that interface wrinkles show spatiotemporal chaos with two length scales. It supplies no rigorous definition, parameter theorem, domain or initial/boundary conditions for biscale chaos. This material is not the antecedent of Sandstede's branch-point question and is not settled by the heat-kernel diagnostic.
\subsection{Short English statement}
The main question asks how branch-point Evans phenomena control decay beyond scalar heat equations. The aid gives a cancellation-sensitive resolvent inversion rule and Neumann/Dirichlet examples with different local powers. Analytic hypotheses, data projections and norms remain essential. The autocatalytic biscale-chaos paragraph belongs to a separate presenter.
\subsection{Sources}
\begin{itemize}
\item Local UnsolvedMath registry, permanent ID 20002445, source code catalogue entry; canonical \texttt{aim-pdes-notes.json}, notes[56].
\item AIM \emph{Stability criteria for multi-dimensional waves and patterns}, B.~Sandstede branch-point question and separate S.~Malham setup: \url{https://aimath.org/WWN/multidimwaves/multidimwaves.pdf}. The supplied source recovery names Murata but gives no full reference for his decay result.
\item ``Biscale chaos in propagating fronts'', \emph{Phys. Rev. E} 52 (1995), pp.~4724--4735, as supplied; authors are not given in this record.
\item Registry research attempt 1, dated 2026-08-14: partial result, rigor ``proved'', novelty confidence low. The explicitly machine-generated, non-peer-reviewed aid requires independent verification before citation.
\end{itemize}
\end{document}
